% Part: normal-modal-logic
% Chapter: axioms-systems
% Section: consistency

\documentclass[../../../include/open-logic-section]{subfiles}

\begin{document}

\olfileid[pt]{nml}{prf}{con}

\olsection{Consistência}

A consistência é uma propriedade importante dos conjuntos de
!!{formula}s. Um conjunto de !!{formula}s é inconsistente se uma
contradição, como~$\lfalse$, é !!{derivable} a partir dele; caso
contrário, é consistente. Se um conjunto é inconsistente, suas
!!{formula}s não podem ser todas verdadeiras em um mundo de um modelo.
Para o teorema da completude, demonstramos a recíproca: todo conjunto
consistente é verdadeiro em um mundo de um modelo, a saber, no
``modelo canônico''.

\begin{defn}
  Um conjunto $\Gamma$ é \emph{consistente} relativamente a um
  sistema~$\Sigma$, ou, como diremos, $\Sigma$-consistente, se e
  somente se $\Gamma
  \Proves/[\Sigma] \lfalse$.
\end{defn}

Assim, por exemplo, o conjunto $\{ \Box(p \lif q), \Box p,
\lnot\Box q \}$ é consistente relativamente à lógica proposicional,
mas não é \Log{K}-consistente. De modo semelhante, o conjunto
$\{ \Diamond p, \Box\Diamond p \lif q, \lnot q \}$ não é
\Log{K5}-consistente.

\begin{prop}\ollabel{prop:consistencyfacts}
  Seja $\Gamma$ um conjunto de !!{formula}s. Então:
  \begin{enumerate}
  \item $\Gamma$ é $\Sigma$-consistente se e somente se existe
    alguma !!{formula}~$!A$ tal que $\Gamma \Proves/[\Sigma]
    !A$.
  \item \ollabel{prop:consistencyfacts-b}%
    $\Gamma \Proves[\Sigma] !A$ se e somente se $\Gamma \cup \{
    \lnot!A \}$ não é $\Sigma$-consistente.
  \item \ollabel{prop:consistencyfacts-c}%
    Se $\Gamma$ é $\Sigma$-consistente, então, para qualquer
    !!{formula} $!A$, ou $\Gamma \cup \{ !A \}$ é
    $\Sigma$-consistente, ou $\Gamma \cup \{ \lnot!A \}$ é
    $\Sigma$-consistente.
  \end{enumerate}
\end{prop}

\begin{proof}
  Esses fatos decorrem facilmente da lógica proposicional clássica.
  Apresentamos o argumento para \olref{prop:consistencyfacts-c}.
  Proceda por contraposição e suponha que nem $\Gamma \cup \{ !A \}$
  nem $\Gamma \cup \{ \lnot!A \}$ seja $\Sigma$-consistente. Então,
  por \olref{prop:consistencyfacts-b}, tanto
  $\Gamma, !A \Proves[\Sigma] \lfalse$ quanto
  $\Gamma, \lnot !A \Proves[\Sigma] \lfalse$. Pelo teorema da dedução,
  $\Gamma \Proves[\Sigma] !A \to \lfalse$ e
  $\Gamma \Proves[\Sigma] \lnot!A \lif \lfalse$. Mas $(!A \lif
  \lfalse) \lif ((\lnot!A \lif \lfalse) \lif \lfalse)$ é uma
  instância tautológica; portanto, por
  \olref[prp]{prop:derivabilityfacts}\olref[prp]{prop:derivabilityfacts-ruleT},
  $\Gamma \Proves[\Sigma] \lfalse$.
\end{proof}

\end{document}
